% =============================================================================
% 5.7 INTERNACIONALIZACIÓN, AUTENTICACIÓN OAUTH Y OPTIMIZACIÓN
% =============================================================================
\section{Internacionalización, Autenticación Federada y Optimización}
\label{sec:i18n_build}

Esta sección documenta tres aspectos transversales del frontend que completan el diseño
detallado: la internacionalización multilingüe, el flujo de autenticación federada OAuth
en el cliente, y las estrategias de optimización del rendimiento de carga.

\subsection{Internacionalización con i18next (es / en / pt)}
\label{subsec:i18n}

La aplicación es \textbf{trilingüe} —español, inglés y portugués— mediante
\textbf{i18next} y \textbf{react-i18next}. Las traducciones se organizan en \textbf{11
espacios de nombres} (\textit{namespaces}) por idioma, lo que segmenta los textos por
área funcional y evita un único fichero monolítico. La tabla~\ref{tab:i18n-ns} enumera los
\textit{namespaces}.

\clearpage
\vspace{0.2cm}
\noindent
\renewcommand{\arraystretch}{1.3}
\fontsize{8.5}{10.2}\selectfont\sffamily
\begin{tabularx}{\textwidth}{|L{3.0cm}|Y|}
\hline
\rowcolor{headerTabla}
\sffamily\textbf{Namespace} & \sffamily\textbf{Contenido} \\
\hline
\texttt{common} · \texttt{nav} & Textos transversales y de navegación. \\ \hline
\rowcolor{grisTecnico}
\texttt{auth} · \texttt{settings} & Pantallas de acceso y de preferencias. \\ \hline
\texttt{dashboard} · \texttt{recruiter} · \texttt{admin} & Espacios de trabajo por rol. \\ \hline
\rowcolor{grisTecnico}
\texttt{public} · \texttt{pages} & Vistas públicas y páginas estáticas. \\ \hline
\texttt{domain} · \texttt{components} & Terminología de dominio y componentes reutilizables. \\ \hline
\end{tabularx}
\label{tab:i18n-ns}

El idioma se resuelve con \comando{i18next-browser-languagedetector}, que infiere la
preferencia del navegador, y se sincroniza con el \comando{authStore} para persistir la
elección del usuario entre sesiones (hidratación de idioma, sección~\ref{subsec:stores},
pág.~\pageref{subsec:stores}). El backend respeta esta preferencia para localizar correos
y mensajes vía \pathbreak{/api/v1/preferences}.

Los componentes consumen las traducciones con el \textit{hook} \comando{useTranslation},
referenciando el \textit{namespace} y la clave, como ilustra el
listado~\ref{lst:i18n-uso}. La inicialización (listado~\ref{lst:i18n-init}) carga los 33
módulos de traducción (11 \textit{namespaces} $\times$ 3 idiomas) y registra el detector.

\begin{lstlisting}[language=JavaScript, caption={Inicialización de i18next con detector de idioma (\texttt{i18n/index.ts}).}, label={lst:i18n-init}]
i18n
  .use(LanguageDetector)
  .use(initReactI18next)
  .init({
    resources: { es, en, pt },     // 11 namespaces por idioma
    fallbackLng: 'es',
    interpolation: { escapeValue: false },
  });
\end{lstlisting}

\begin{lstlisting}[language=JavaScript, caption={Consumo de traducciones en un componente.}, label={lst:i18n-uso}]
const { t } = useTranslation('dashboard');
return <h1>{t('portfolio.title')}</h1>;  // -> "Mi portafolio" / "My portfolio"
\end{lstlisting}

\subsection{Flujo de autenticación federada (OAuth2) en el cliente}
\label{subsec:oauth}

Además del acceso con credenciales, EthosHub admite \textbf{OAuth2} (proveedores sociales
vía Supabase Auth). El cliente no maneja secretos de OAuth: delega el intercambio en
Supabase y solo recibe el resultado en una ruta de \textit{callback}. La
figura~\ref{fig:flow-oauth} ilustra el flujo.

\vspace{0.3cm}
\begin{figure}[h!]
\centering
\begin{tikzpicture}[node distance=0.7cm and 1.3cm,
    box/.style={rectangle, rounded corners=2pt, draw, font=\sffamily\scriptsize, align=center,
        minimum width=2.6cm, minimum height=0.85cm, inner sep=3pt},
    rel/.style={-{Stealth[length=2mm]}, semithick}]
    \node[box, fill=c4Persona!25] (user) {Usuario};
    \node[box, fill=c4Componente!40, right=of user] (login) {LoginPage\\\scriptsize ``Entrar con...''};
    \node[box, fill=colorAcento!28, draw=colorAcento!80!black, right=of login] (sb) {Supabase\\Auth (OAuth)};
    \node[box, fill=c4Componente!40, below=of sb] (cb) {OAuth2Callback\\Page};
    \node[box, fill=c4Componente!45, below=of login] (sync) {authService\\\scriptsize POST /oauth/sync};
    \node[box, fill=c4Sistema, text=white, left=of sync] (home) {Dashboard\\por rol};

    \draw[rel] (user) -- (login);
    \draw[rel] (login) -- node[c4etiqueta, above]{\scriptsize redirect} (sb);
    \draw[rel] (sb) -- node[c4etiqueta, right]{\scriptsize token} (cb);
    \draw[rel] (cb) -- (sync);
    \draw[rel] (sync) -- node[c4etiqueta, above]{\scriptsize sesión} (home);
\end{tikzpicture}
\caption{Flujo de autenticación federada OAuth2 en el cliente.}
\label{fig:flow-oauth}
\end{figure}

Tras el retorno, \comando{OAuth2CallbackPage} extrae el token de Supabase y llama a
\comando{authService} para sincronizar el perfil en el backend
(\pathbreak{POST /api/auth/oauth/sync}, capítulo~\ref{ch:backend}, tabla~\ref{tab:endpoints},
pág.~\pageref{tab:endpoints}), creando el perfil local si es la primera vez. El servidor
de desarrollo de Vite establece la cabecera \\
\comando{Cross-Origin-Opener-Policy: same-origin-allow-popups} para permitir el flujo de
ventana emergente de OAuth.

\subsection{Optimización del rendimiento de carga}
\label{subsec:optimizacion}

El rendimiento percibido se optimiza con tres técnicas combinadas, descritas en la
tabla~\ref{tab:optimizacion}.

\vspace{0.2cm}
\noindent
\renewcommand{\arraystretch}{1.3}
\fontsize{8.5}{10.2}\selectfont\sffamily
\begin{tabularx}{\textwidth}{|L{3.6cm}|Y|}
\hline
\rowcolor{headerTabla}
\sffamily\textbf{Técnica} & \sffamily\textbf{Efecto} \\
\hline
\textit{Code splitting} por ruta & Cada página se carga bajo demanda (\texttt{React.lazy} + \texttt{Suspense}). \\ \hline
\rowcolor{grisTecnico}
\textit{Manual chunks} (Vite) & Se aísla el \textit{vendor} (\texttt{react}, \texttt{react-dom}, \texttt{react-router-dom}) en un \textit{chunk} cacheable. \\ \hline
\textit{Optimize deps} & Se pre-empaquetan \texttt{leaflet} y \texttt{react-leaflet} para acelerar el arranque en desarrollo. \\ \hline
\rowcolor{grisTecnico}
\textit{Skeletons} de carga & \texttt{Suspense} muestra esqueletos en vez de pantallas en blanco durante la carga diferida. \\ \hline
\end{tabularx}
\label{tab:optimizacion}

\begin{NotaTecnica}
La división del \textit{vendor} en su propio \textit{chunk} es una optimización de caché:
las dependencias base (React y el enrutador) cambian con poca frecuencia, por lo que el
navegador puede reutilizar ese \textit{chunk} entre despliegues mientras solo se
re-descarga el código de aplicación modificado.
\end{NotaTecnica}